%%%PAGE 132 rot=0 legibility=good summary=天体密度公式；重要(第一宇宙速度)、合并(Tρ关系)、表面(g∝ρR)三个结论；先表达后正比的比例技巧；两极与赤道测密度
%%%BLOCK id=132-01 ch=05 sec=天体质量密度 kind=formula alt=none cont=prev y=0.08-0.31
\textbf{4. 密度}
\[ \rho=\frac{M}{\frac{4}{3}\pi R^{3}} \]

1. 重要\quad $V_1=\sqrt{\dfrac{gM}{R}}=\sqrt{gR}$ % CHECK: 原稿根号内分子字母形似小写 g（物理上应为 G）

\[
\text{推导}\
\left\{\begin{aligned}
\frac{GMm}{R^{2}}&=\frac{mv^{2}}{R}\ \text{(行星)}\\
\frac{GMm'}{R^{2}}&=m'g\ \text{(表面)}
\end{aligned}\right.
\qquad \therefore\ mg=\frac{mv^{2}}{R}\qquad V=\sqrt{gR}
\]
%%%END
%%%BLOCK id=132-02 ch=05 sec=天体质量密度 kind=formula alt=none cont=none y=0.31-0.55
2. 合并
\[
\left\{\begin{aligned}
\frac{GMm}{R^{2}}&=\frac{m\,4\pi^{2}R}{T^{2}}\\
\rho&=\frac{M}{\frac{4}{3}\pi R^{3}}
\end{aligned}\right.
\qquad
T=\sqrt{\frac{4\pi R^{3}}{GM}}=\sqrt{\frac{4\pi^{2}R^{3}}{G\rho\,\frac{4}{3}\pi R^{3}}}=\sqrt{\frac{3\pi}{G\rho}}
\] % CHECK: 第一个根号内分子原稿为 4πR^3（缺 π 的平方，应为 4π^2R^3）

有3没4\quad $\rho$\,1\;$T$\,2\;(次方)

3. 表面
\[
g=\frac{GM}{R^{2}}=\frac{G\rho\,\frac{4}{3}\pi R^{3}}{R^{2}}=\frac{4}{3}\pi G\rho R\propto\rho R
\]
\[ \rho=\frac{M}{\frac{4}{3}\pi R^{3}} \]
%%%END
%%%BLOCK id=132-03 ch=05 sec=比例法·先表达后正比 kind=method alt=90 cont=none y=0.55-0.76
4. 先表达后正比\quad 常量都当 1\ $\propto$
\[
\begin{array}{cl}
\text{\unclear{O}} & \propto\ \dfrac{a}{b}\\ \cline{1-1}
\text{\unclear{O}} & \propto\ \dfrac{a'}{b'}
\end{array}
\qquad
\frac{a}{b}\div\frac{a'}{b'}=\frac{a\div a'}{b\div b'}
\qquad
\frac{\dfrac{b}{a}}{\dfrac{d}{c}}=\frac{bc}{ad}
\]
\[
\frac{\dfrac{\text{①}}{\text{②}}}{\dfrac{\text{③}}{\text{④}}}=\frac{\text{①④}}{\text{②③}}
\qquad \text{中沉两头合}
\]
%%%END
%%%BLOCK id=132-04 ch=05 sec=天体质量密度·两极赤道 kind=derivation alt=none cont=none y=0.76-0.84
\[
\begin{aligned}
\text{两极}\quad mg_0&=\frac{GMm}{R^{2}}=\frac{4}{3}\pi\rho mGR\\
\text{赤\struck{极}道}\quad mg+\frac{m\,4\pi^{2}R}{T^{2}}&=\frac{GMm}{R^{2}}=\frac{4}{3}\pi\rho mGR\qquad \rho=\frac{3\pi}{GT^{2}}\cdot\frac{g_0}{g_0-g}
\end{aligned}
\] % 原稿“赤”后多写了一个被涂改的“极”字，“道”字换行写在下方
% 页面左侧边缘 y≈0.77 处另有一小段红笔痕迹（疑为背面透字或涂鸦），未转录
%%%END
%%%PAGEEND 132
